\subsection{Affine Reduction and Event-Conditioned Controls}

The first step tested whether transition-linked motion could be explained by
affine geometry. Whole-swarm affine motion captures translation and deformation
at the global scale. Local affine motion captures the best local deformation of
the focal neighborhood. T1 was computed after the local affine component was
subtracted, so a positive T1 result indicates residual motion that is not fully
absorbed by the local affine fit.

For each observation, event-conditioned T1 values were compared with
non-event or shifted-event controls. Robustness was then checked across nearby
local-neighborhood scales and temporal lags. This produced two kinds of
evidence: whether T1 survived at the original settings, and whether that
survival remained stable under local scale/lag perturbations.

\subsection{Spatial and Timing Decomposition}

To interpret the form of the residual, we separated diffuse tangential
activity from more localized or directional alternatives. The main spatial
question was whether all-tangential activity was more repeatable than an
edge/core contrast. The timing question was whether activity was concentrated
in the near-pre event window. Signed event-type analyses were treated as
boundary tests because their cross-observation direction and order were not
stable enough to support a universal rule.

